\documentclass[11pt,a4paper]{article}

\usepackage[utf8]{inputenc}
\usepackage[T1]{fontenc}
\usepackage[margin=1in,headheight=14pt]{geometry}
\usepackage{amsmath,amssymb}
\usepackage{xcolor}
\usepackage{tabularx}
\usepackage{booktabs}
\usepackage{hyperref}
\usepackage{fancyhdr}
\usepackage{tcolorbox}
\usepackage{enumitem}
\usepackage{listings}

% Palette Definition
\definecolor{primaryblue}{RGB}{24, 76, 120}
\definecolor{accentred}{RGB}{180, 40, 40}
\definecolor{accentgreen}{RGB}{34, 139, 34}
\definecolor{accentamber}{RGB}{204, 102, 0}
\definecolor{darkslate}{RGB}{30, 41, 59}
\definecolor{lightbg}{RGB}{248, 250, 252}
\definecolor{codebg}{RGB}{241, 245, 249}

\hypersetup{
    colorlinks=true,
    linkcolor=primaryblue,
    citecolor=primaryblue,
    urlcolor=primaryblue
}

\tcbuselibrary{skins,breakable}

\newtcolorbox{alertbox}[2][]{
    enhanced,
    colback=red!5!white,
    colframe=accentred,
    fonttitle=\bfseries\color{white},
    colbacktitle=accentred,
    title=#2,
    arc=4mm,
    boxrule=1pt,
    left=6pt, right=6pt, top=6pt, bottom=6pt,
    #1
}

\newtcolorbox{infobox}[2][]{
    enhanced,
    colback=blue!5!white,
    colframe=primaryblue,
    fonttitle=\bfseries\color{white},
    colbacktitle=primaryblue,
    title=#2,
    arc=4mm,
    boxrule=1pt,
    left=6pt, right=6pt, top=6pt, bottom=6pt,
    #1
}

\newtcolorbox{conflictbox}[2][]{
    enhanced,
    colback=orange!6!white,
    colframe=accentamber,
    fonttitle=\bfseries\color{white},
    colbacktitle=accentamber,
    title=#2,
    arc=4mm,
    boxrule=1.2pt,
    left=6pt, right=6pt, top=6pt, bottom=6pt,
    #1
}

\lstset{
    backgroundcolor=\color{codebg},
    basicstyle=\ttfamily\small\color{darkslate},
    keywordstyle=\bfseries\color{primaryblue},
    commentstyle=\itshape\color{gray},
    stringstyle=\color{accentred},
    breaklines=true,
    frame=single,
    rulecolor=\color{gray!30},
    numbers=left,
    numberstyle=\tiny\color{gray},
    numbersep=8pt,
    tabsize=2,
    showstringspaces=false
}

\pagestyle{fancy}
\fancyhf{}
\fancyhead[L]{\textsc{Test DPC Adult Content Defense Engine}}
\fancyhead[R]{\thepage}
\fancyfoot[C]{\footnotesize Technical Specification \& Conflict Analysis}
\renewcommand{\headrulewidth}{0.4pt}
\renewcommand{\footrulewidth}{0.4pt}

\title{
    \vspace{-1.5cm}
    \Huge\textbf{\color{primaryblue}Test DPC Adult Content Defense Engine}\\[0.4cm]
    \Large\textbf{Comprehensive Architecture, Logic Specification, and Subsystem Conflict Analysis}\\[0.2cm]
    \normalsize An In-Depth Technical Specification of Android Enterprise Device Owner Enforcements, Multi-Model LLM Heuristics, Real-Time Accessibility Traps, and Architectural Remediation
}
\author{\textbf{DpcLocker Architecture Team}}
\date{August 16, 2026}

\begin{document}

\maketitle

\begin{abstract}
This document provides a formal, comprehensive architectural specification of the adult content blocking, anti-impulse, and perimeter defense systems integrated into the customized Android Enterprise Device Policy Controller (\texttt{TestDPC} / \texttt{DpcLocker}). We analyze the deterministic, heuristic, and generative-AI layers spanning runtime search and screen interception (\texttt{ImpulseGuardService}, \texttt{GeminiGuardEngine}), progressive escalating disciplinary penalties (\texttt{PenaltyManager}), proactive installation traps (\texttt{BrowserBlocker}, \texttt{NotoriousAppBlocker}, \texttt{AiAppAuditor}), daily foreground time limits (\texttt{AppTimerManager}), and Google Chrome Enterprise Managed Configurations (\texttt{ChromePolicyManager}). Finally, we present an exhaustive audit of five critical structural conflicts currently creating system instability, race conditions, and accidental un-suspension bugs across the active codebase, concluding with a mathematically rigorous refactoring blueprint.
\end{abstract}

\tableofcontents
\newpage

\section{System Architecture Overview}

The adult content defense system operates as a five-layer defense-in-depth perimeter deployed inside an Android Enterprise Device Owner application. The primary design goal is \textbf{impulse interruption}---the elimination of friction-free access to adult content, unmanaged third-party web browsers, video scrapers, and bypass mechanisms.

\begin{table}[h!]
\centering
\small
\begin{tabularx}{\textwidth}{l p{3.2cm} p{4.2cm} X}
\toprule
\textbf{Layer} & \textbf{Component} & \textbf{Mechanism} & \textbf{Target Threat Model} \\
\midrule
\textbf{Layer 1} & \texttt{ImpulseGuardService} \newline \texttt{GeminiGuardEngine} & Real-time Accessibility Node Parsing, 0ms Disk/RAM Cache, Multi-Model Flash-Lite Ladder & Explicit search queries, adult keywords, and on-screen textual erotic content. \\
\addlinespace
\textbf{Layer 2} & \texttt{PenaltyManager} & Exponential Daily Violation Ladder ($1\text{m} \to 10\text{m} \to 30\text{m}$), 24h Reset Timer & Habitual compulsive browsing, repetitive bypass attempts. \\
\addlinespace
\textbf{Layer 3} & \texttt{BrowserBlocker} \newline \texttt{NotoriousAppBlocker} \newline \texttt{AiAppAuditor} & Intent Filter Scans, Static Social Package Blocks, Manifest \& Permission AI Audits & Installation of alternative browsers (Firefox, Brave), video downloaders (Snaptube, Vmate), and social feeds (X, Reddit). \\
\addlinespace
\textbf{Layer 4} & \texttt{AppTimerManager} & \texttt{UsageStatsManager} API, Foreground Time Quotas, Reflection Observers & Screen addiction, extended late-night usage of potentially vulnerable apps. \\
\addlinespace
\textbf{Layer 5} & \texttt{ChromePolicyManager} & Android Enterprise Managed Configurations (\texttt{dpm.setApplicationRestrictions}) & Google SafeSearch bypass, Proxy/VPN extensions, Incognito mode, explicit domains (\texttt{fboxtv.org}). \\
\bottomrule
\end{tabularx}
\caption{The Five Integrated Defense Layers of the Test DPC Security Core.}
\label{tab:layers}
\end{table}

---

\section{Layer 1: Real-Time Runtime Content Filtering}

Layer 1 operates as a privileged background system inspection service built upon the Android \texttt{AccessibilityService} framework (\texttt{ImpulseGuardService.java}) coupled with a low-latency multi-model Gemini API inference pipeline (\texttt{GeminiGuardEngine.java}).

\subsection{Event Interception \& Filtering Pipeline}
The service subscribes to core OS window and view events:
\begin{itemize}[leftmargin=1.5cm]
    \item \textbf{\texttt{TYPE\_VIEW\_TEXT\_CHANGED}:} Captures live character entry in edit text nodes and URL address bars.
    \item \textbf{\texttt{TYPE\_VIEW\_CLICKED}:} Captures direct user execution (e.g., clicking a ``Search'' magnifying glass icon or a search recommendation).
    \item \textbf{\texttt{TYPE\_WINDOW\_CONTENT\_CHANGED} \& \texttt{TYPE\_VIEW\_SCROLLED}:} Intercepts rendering events to trigger asynchronous full-screen textual inspection.
\end{itemize}

\subsection{Privacy Masking \& Whitelisting Logic}
To prevent performance degradation and maintain strict user privacy:
\begin{enumerate}
    \item \textbf{Password Field Guard:} If \texttt{event.isPassword()} evaluates to \texttt{true}, or if any parent/source accessibility node contains the substring \texttt{"password"}, the extraction is immediately aborted.
    \item \textbf{System App Whitelist:} Essential system components (\texttt{com.android.systemui}, \texttt{com.google.android.inputmethod.latin}, \texttt{com.android.settings}, \texttt{com.whatsapp}, \texttt{com.google.android.gm}) bypass evaluation instantly with 0ms overhead.
    \item \textbf{URL \& Syntax Filter:} Text strings starting with \texttt{http://}, \texttt{https://}, \texttt{www.}, or generic browser placeholder strings (e.g., \texttt{"Search or type URL"}) are discarded.
\end{enumerate}

\subsection{Concurrency Control \& Debouncing Model}
Because continuous keystroke evaluation causes network throttling and server load, \texttt{ImpulseGuardService} implements a hybrid dispatch mechanism:
\begin{itemize}
    \item \textbf{Fast-Path Immediate Dispatch:} If an explicit click event (\texttt{TYPE\_VIEW\_CLICKED}) is detected, the query is dispatched instantly.
    \item \textbf{1200ms Debounce Pause:} For typing events, a trailing-edge debounce runnable delays dispatch until the user pauses for at least $1200\,\text{ms}$, ensuring that the complete intended query is evaluated rather than partial keystroke prefixes.
    \item \textbf{Single-Flight Concurrency Lock:} An \texttt{AtomicBoolean} (\texttt{isRequestInFlight}) prevents multiple overlapping network threads. Intermediate keystrokes during an active inference call are safely dropped.
\end{itemize}

\subsection{0ms Two-Tier Cache Architecture}
Before any outbound HTTPS request is generated, queries are checked against a dual-tier cache:
\begin{equation}
\text{Verdict}(q) = 
\begin{cases} 
\text{Hit}(\text{Risky}), & \text{if } q \in \text{DiskCache}_{\text{risky}} \implies \text{Instant Suspension} \\
\text{Hit}(\text{Safe}), & \text{if } q \in \text{DiskCache}_{\text{safe}} \implies \text{Instant Pass} \\
\text{Miss}, & \implies \text{Gemini Multi-Model Fallback Ladder}
\end{cases}
\end{equation}

For on-screen plain text, a 500-entry \texttt{LruCache<String, Boolean>} in RAM stores deterministic hash keys ($H = \text{hash}(\text{screenText})$). To prevent false-negative persistence on dynamic pages, \textbf{only risky screen verdicts} are persisted in RAM and disk.

\subsection{Generative AI Fallback Ladder}
When a cache miss occurs, the engine executes an HTTP POST against Google's Generative Language API using a resilient fallback model hierarchy:
\begin{equation}
\mathcal{M} = \big[ \text{3.5-flash-lite} \longrightarrow \text{3.1-flash-lite} \longrightarrow \text{2.5-flash-lite} \longrightarrow \text{1.5-flash-8b} \big]
\end{equation}
Each model is queried sequentially with a strict $8000\,\text{ms}$ socket timeout. If HTTP 200 is received, a fuzzy JSON parser extracts the boolean flag \texttt{"is\_risky": true/false}. A fail-safe open policy applies if all models fail or the network is unreachable.

---

\section{Layer 2: Progressive Penalty Ladder \& Suspension Engine}

When Layer 1 identifies explicit adult intent, the target application is not merely closed; it is placed into OS-level package quarantine via the Android Enterprise Device Administration API.

\subsection{Mathematical Model of Progressive Escalation}
The severity of the suspension duration increases monotonically with each violation recorded within a sliding 24-hour window:
\begin{equation}
D(v) = 
\begin{cases}
P_1 = 1\,\text{minute}, & v = 1 \\
P_2 = 10\,\text{minutes}, & v = 2 \\
P_3 = 30\,\text{minutes}, & v = 3 \\
P_4 = 30\,\text{minutes}, & v \ge 4
\end{cases}
\end{equation}
where $v$ is the cumulative daily violation count, and $P_1, P_2, P_3, P_4$ are user-configurable in the Test DPC settings screen.

\subsection{The 24-Hour Inactivity Reset Rule}
Let $t_{\text{current}}$ be the timestamp of the current event, and $t_{\text{last}}$ be the timestamp of the previous recorded violation. The violation counter resets according to the step function:
\begin{equation}
v_{\text{new}} = 
\begin{cases}
1, & \text{if } (t_{\text{current}} - t_{\text{last}}) \ge 86{,}400{,}000\,\text{ms} \quad (24\text{ hours}) \\
v_{\text{prev}} + 1, & \text{otherwise}
\end{cases}
\end{equation}

\subsection{Operating System Enforcement Mechanism}
Suspension is executed via privileged IPC to the Android \texttt{DevicePolicyManagerService}:
\begin{lstlisting}[language=Java]
DevicePolicyManager dpm = (DevicePolicyManager) context.getSystemService(Context.DEVICE_POLICY_SERVICE);
ComponentName admin = DeviceAdminReceiver.getComponentName(context);
dpm.setPackagesSuspended(admin, new String[]{packageName}, true);
\end{lstlisting}
When suspended:
\begin{itemize}
    \item The target application's launcher icon is instantly greyed out.
    \item Any active foreground activity is forcefully killed.
    \item Launch intents from other applications or notifications are intercepted by the OS with a dialog indicating the app has been suspended by the device administrator.
\end{itemize}

---

\section{Layer 3: Install-Time Defense \& Proactive Sandboxing}

To prevent users from circumventing Chrome filters by installing alternative web browsers or video downloading tools from APK files or secondary app stores, Layer 3 enforces install-time inspection.

\subsection{Browser Neutralization Engine (\texttt{BrowserBlocker})}
\texttt{BrowserBlocker} listens for newly added packages via \texttt{LauncherApps.Callback} and broadcast intents. It dynamically tests whether the new package declares support for generic, browsable web intents:
\begin{lstlisting}[language=Java]
Intent webIntent = new Intent(Intent.ACTION_VIEW, Uri.parse("https://www.google.com"));
webIntent.addCategory(Intent.CATEGORY_BROWSABLE);
List<ResolveInfo> handlers = pm.queryIntentActivities(webIntent, PackageManager.MATCH_ALL);
\end{lstlisting}
If the package resolves HTTP/HTTPS viewing and is not in the system whitelist or Google Chrome (\texttt{com.android.chrome}), it is \textbf{permanently suspended} immediately upon installation.

\subsection{Social Media \& Content Blocklist (\texttt{NotoriousAppBlocker})}
A persistent blacklist targets notorious platforms known for uncurated adult content feeds:
\begin{itemize}
    \item X / Twitter (\texttt{com.twitter.android}, \texttt{com.twitter.android.lite})
    \item Reddit (\texttt{com.reddit.frontpage})
    \item Tumblr (\texttt{com.tumblr})
    \item Telegram (\texttt{org.telegram.messenger}, \texttt{org.telegram.messenger.web}, \texttt{org.telegram.plus})
    \item Stock YouTube (\texttt{com.google.android.youtube}, enforcing clean ReVanced builds)
\end{itemize}

\subsection{Deep Manifest AI Auditor (\texttt{AiAppAuditor})}
For unknown non-browser apps, \texttt{AiAppAuditor} extracts the full Android manifest profile:
\begin{enumerate}
    \item Package Name, Human-Readable Label, and Application Category (\texttt{ApplicationInfo.category}).
    \item All requested permissions (specifically \texttt{INTERNET}, storage permissions, foreground service).
    \item Declared Activity names (searching for \texttt{WebView}, \texttt{Downloader}, \texttt{Player}, \texttt{MediaFetch}).
\end{enumerate}
This structured JSON is sent to Gemini Flash models with a strict system instruction to classify hidden video downloaders (e.g., Snaptube, VidMate, XNXX apps). If the AI is offline, a deterministic structural fallback flags any app possessing $\text{INTERNET} \land \text{STORAGE} \land \text{WebView/Downloader Components}$.

---

\section{Layer 4: Managed Chrome Restrictions \& Enterprise Policies}

Layer 4 leverages Android Enterprise Managed Configurations (\texttt{ChromePolicyManager.java}) to enforce hardware-level policy control inside Google Chrome:

\begin{lstlisting}[language=Java]
Bundle chromeBundle = new Bundle();
// Force Google SafeSearch
chromeBundle.putBoolean("ForceGoogleSafeSearch", true);
chromeBundle.putInt("SafeSearchMode", 1);

// Enable Chrome SafeSites Adult Content Filter
chromeBundle.putInt("SafeSitesFilterBehavior", 1);

// Force Direct Connection (Disables Proxy & VPN Bypass extensions)
chromeBundle.putString("ProxyMode", "direct");

// Enforce Domain Filtering Blocklist
String[] urlBlocklist = new String[]{
    "fboxtv.org", "x.com", "twitter.com", "reddit.com",
    "tumblr.com", "croxyproxy.com", "proxysite.com",
    "youtube.com/shorts/*", "*://*.youtube.com/shorts/*"
};
chromeBundle.putStringArray("URLBlocklist", urlBlocklist);
chromeBundle.putStringArray("URLBlacklist", urlBlocklist);

dpm.setApplicationRestrictions(admin, "com.android.chrome", chromeBundle);
\end{lstlisting}

---

\section{Critical Subsystem Conflicts \& Vulnerability Analysis}

Our exhaustive line-by-line audit identified \textbf{five severe architectural conflicts} that cause system instability, race conditions, and un-suspension leaks.

\begin{conflictbox}{Conflict 1: Temporary vs. Permanent Suspension Overwrite (The Unsuspend Leak)}
\textbf{Root Cause:} \texttt{ImpulseGuardService.java:185} vs. \texttt{BrowserBlocker.java:140} \& \texttt{NotoriousAppBlocker.java:81}.\\[0.2cm]
\textbf{The Mechanism:}
\begin{enumerate}
    \item \texttt{BrowserBlocker} or \texttt{NotoriousAppBlocker} permanently suspends an app (e.g., Firefox or Twitter) upon installation.
    \item If that application is also listed in \texttt{ImpulseGuardService}'s monitored list and triggers an Accessibility event or lingering audit, \texttt{ImpulseGuardService} records a timed suspension entry in \texttt{PREF\_SUSPENSIONS} (e.g., $1\,\text{minute}$).
    \item When the $1\,\text{minute}$ timer elapses, \texttt{checkAndCleanExpiredSuspensions()} executes:
    $$\texttt{dpm.setPackagesSuspended(admin, new String[]\{pkg\}, false);}$$
    \item \textbf{Critical Failure:} The permanently blocked web browser or social app is \textbf{un-suspended and unlocked for the user}, completely breaking the perimeter defense!
\end{enumerate}
\end{conflictbox}

\begin{conflictbox}{Conflict 2: Triple-Receiver Installation Stampede}
\textbf{Root Cause:} Redundant declarations in \texttt{AndroidManifest.xml}.\\[0.2cm]
When an APK is installed, three separate listeners fire in parallel:
\begin{enumerate}
    \item \texttt{PackageInstallReceiver.java} (BroadcastReceiver on \texttt{ACTION\_PACKAGE\_ADDED})
    \item \texttt{PackageInstallationReceiver.java} (BroadcastReceiver on \texttt{ACTION\_PACKAGE\_ADDED})
    \item \texttt{LauncherApps.Callback.onPackageAdded()} inside \texttt{BrowserBlocker.java}
\end{enumerate}
This results in redundant simultaneous calls to \texttt{BrowserBlocker.checkAndSuspendPackage()} and creates concurrent write locks on \texttt{SharedPreferences}.
\end{conflictbox}

\begin{conflictbox}{Conflict 3: Subsystem Whitelist Contradictions}
\textbf{Root Cause:} \texttt{AiAppAuditor.java:109} vs. \texttt{NotoriousAppBlocker.java:23}.\\[0.2cm]
\texttt{AiAppAuditor} includes \texttt{"org.telegram.messenger"} in its hardcoded trusted bypass list to prevent video downloader scans, while \texttt{NotoriousAppBlocker} explicitly marks \texttt{"org.telegram.messenger"} as a default blocked package. Depending on which receiver executes first, logs and states emit conflicting telemetry.
\end{conflictbox}

\begin{conflictbox}{Conflict 4: Fragmented API Key Storage}
\textbf{Root Cause:} \texttt{GeminiGuardEngine.PREFS\_NAME} (\texttt{"gemini\_guard\_engine\_prefs"}) vs. \texttt{AiAppAuditor.PREFS\_NAME} (\texttt{"ai\_app\_auditor\_prefs"}).\\[0.2cm]
The runtime keystroke guard and the install-time manifest auditor read from two completely isolated storage files. If the user configures their Gemini API key in the Impulse Guard dialog, the install-time AI Auditor remains disabled and falls back to coarse heuristics.
\end{conflictbox}

\begin{conflictbox}{Conflict 5: Chrome Incognito Availability Mismatch}
\textbf{Root Cause:} \texttt{ChromePolicyManager.java:26}.\\[0.2cm]
Line 26 sets \texttt{chromeBundle.putInt("IncognitoModeAvailability", 0);}, which explicitly enables Incognito mode. Standard adult protection requires \texttt{IncognitoModeAvailability = 1} (Disabled) to ensure URL blocklists, SafeSites filtering, and network logging cannot be bypassed in private browsing tabs.
\end{conflictbox}

---

\section{Architectural Reconciliation \& Remediation Blueprint}

To eliminate all five bugs and achieve a robust unified architecture:

\begin{enumerate}
    \item \textbf{Unified Suspension Category Hierarchy:}
    Create an enumerated suspension registry:
    $$\text{State}(\text{pkg}) \in \{\text{UNSUSPENDED},\; \text{TEMPORARY\_PENALTY},\; \text{PERMANENT\_PROHIBITED},\; \text{TIMER\_EXCEEDED}\}$$
    \texttt{ImpulseGuardService.checkAndCleanExpiredSuspensions()} must query \texttt{NotoriousAppBlocker.isPackageBlocked()} and \texttt{BrowserBlocker.isNonChromeBrowser()} prior to calling \texttt{dpm.setPackagesSuspended(false)}. If the app is permanently prohibited, the auto-unsuspend call is suppressed.

    \item \textbf{Receiver Consolidation:}
    Remove \texttt{PackageInstallReceiver.java} and \texttt{PackageInstallationReceiver.java} from \texttt{AndroidManifest.xml}. Route all install events exclusively through \texttt{LauncherApps.Callback} inside a unified \texttt{SecurityPipelineManager}.

    \item \textbf{Single Secrets Vault:}
    Refactor \texttt{AiAppAuditor} and \texttt{GeminiGuardEngine} to share a common key provider:
    $$\texttt{SecurityConfig.getGeminiApiKey(context)}$$

    \item \textbf{Enforce Incognito Lock:}
    Update \texttt{ChromePolicyManager.java} to enforce \texttt{IncognitoModeAvailability = 1}.
\end{enumerate}

\begin{infobox}{Conclusion}
The Test DPC defense engine possesses an industry-leading multi-layered architecture combining generative AI with Android Device Owner APIs. Implementing the reconciliation blueprint eliminates all five identified race conditions, delivering an un-bypassable, self-healing protection system.
\end{infobox}

\end{document}
